\documentclass[10pt,a4paper]{article}
\usepackage[utf8]{inputenc}
\usepackage[T1]{fontenc}
\usepackage{lmodern}
\usepackage[margin=2.3cm]{geometry}
\usepackage{amsmath,amssymb,booktabs,microtype,enumitem}
\usepackage[hidelinks]{hyperref}
\title{\textbf{SotSort v8-r}\\[3pt]\large Adenda al resumen y al libro: reconstrucción,\\
optimizaciones nuevas y planificación fuera de orden}
\author{Lucas Sotomayor}
\date{17 de noviembre de 2026; sección~\ref{sec:r6} agregada el 19 de noviembre --- borrador}
\renewcommand{\tablename}{Tabla}
\renewcommand{\contentsname}{Índice}
\begin{document}
\maketitle

\section{Por qué existe esta adenda}
El código fuente de la versión 8, sus pruebas y los datos de las figuras se perdieron.
\texttt{sotsort.hpp} (``v8-r'') es una reescritura hecha desde el pseudocódigo, las
demostraciones, el libro y una versión anterior del código. Consecuencias:
\begin{itemize}[nosep]
\item \textbf{Los números del resumen y de los capítulos 1 y 2 del libro son del código
  perdido.} No se pueden reproducir y no describen a v8-r. Los vigentes son los de esta
  adenda.
\item Los algoritmos y las demostraciones de esos textos siguen valiendo para el
  pseudocódigo. Lo nuevo está demostrado en las secciones 13 y 14 del documento de
  demostraciones.
\item Las mediciones de acá son de una sola máquina: una VM de 2 vCPU (Xeon con
  AVX-512), gcc 13, $n=10^6$ salvo que se indique, mediana de 5 corridas fijadas a un
  núcleo. Dos corridas idénticas llegaron a diferir en $\pm15\%$. No deberían citarse
  sin repetirlas en hardware dedicado.
\end{itemize}

\section{Qué se agregó}
\begin{description}[leftmargin=1.2em,itemsep=2pt]
\item[Sonda de rachas.] En cada mezcla con buffer del camino de runs naturales, 8
  búsquedas binarias acotan la cantidad de cambios de lado. Con rachas largas se usa
  galope; si no, \emph{parity merge} cuando los dos runs entran en el buffer, o el lazo
  clásico. Solo para tipos que se copian byte a byte.
\item[Sonda de orden global y regla de ruta.] Con pocos giros, 256 muestras en
  posiciones pseudoaleatorias dicen si los runs tienen las claves entremezcladas. Si es
  así, se va al quicksort salvo que haya pocos runs: se compara
  $k\lceil\lg R\rceil$ con $\lceil\lg(n/L)\rceil$, con $k$ entre $5/4$ y $3$ según el
  modo. Corrige el error conocido de mandar ``runs de 63'' al camino de mezclas.
\item[Zonas desordenadas.] Si el camino de runs naturales encuentra 8 o más runs
  seguidos de menos de 8 elementos, ordena toda la zona con el quicksort estable y la
  trata como un run.
\item[Recorridos vectoriales.] La detección de orden y el conteo de giros comparan 8 o
  16 pares por instrucción con enteros.
\item[Selección sin salto en gcc.] En los \emph{parity merges}, gcc convertía
  \texttt{t ? *a : *b} en un salto para estructuras; clang no. Se reemplazó por una
  selección con aritmética de punteros, solo para tipos no aritméticos y solo en gcc.
\item[Planificador experimental.] Sección~\ref{sec:ooo}.
\end{description}
Ninguna adición es un procedimiento de ordenamiento nuevo: todas eligen entre
procedimientos ya demostrados o cambian cómo se recorre.

\section{Efecto medido de las optimizaciones}
La tabla~\ref{tab:ab} compara la reconstrucción inicial con la actual, en el mismo
binario y con corridas intercaladas (mínimo de 7).

\begin{table}[h]\centering\footnotesize
\begin{tabular}{lccccc}
\toprule
Entrada & u64 & u32 & par (struct) & u64 con comparador & \emph{strings}\\
\midrule
aleatoria & 0,99 / 0,97 & 0,98 / 0,97 & 0,81 / 0,76 & 1,02 / 1,00 & 1,02 / 1,02\\
16 valores & 0,96 / 0,98 & 0,97 / 1,01 & 1,01 / 0,97 & 1,00 / 0,97 & 1,03 / 1,01\\
Zipf & 0,99 / 0,93 & 0,95 / 1,01 & 0,99 / 0,91 & 1,01 / 0,95 & 1,01 / 0,94\\
ordenada & 0,86 / 0,83 & 0,44 / 0,45 & 1,01 / 1,01 & 0,99 / 1,00 & 0,80 / 0,81\\
invertida & 0,89 / 0,89 & 0,62 / 0,62 & 0,99 / 0,96 & 1,01 / 1,27 & 1,05 / 1,10\\
1\% de ruido & 0,78 / 0,94 & 0,62 / 0,94 & 1,06 / 1,01 & 0,65 / 0,98 & 0,99 / 0,94\\
cola 10\% al azar & 0,51 / 0,66 & 0,45 / 0,61 & 0,52 / 0,79 & 0,66 / 0,72 & 1,00 / 0,97\\
10\% al azar en el medio & 0,50 / 0,65 & 0,45 / 0,61 & 0,52 / 0,81 & 0,66 / 0,72 & 1,01 / 0,98\\
20 bloques al azar & 0,91 / 0,86 & 0,79 / 0,84 & 0,89 / 1,01 & 0,86 / 0,94 & 1,00 / 0,94\\
\emph{organ pipe} & 1,06 / 0,89 & 0,56 / 0,79 & 0,62 / 0,99 & 0,92 / 0,99 & 1,00 / 0,88\\
sierra de 8 dientes & 1,06 / 0,91 & 0,73 / 0,90 & 0,54 / 1,04 & 0,89 / 0,98 & 1,00 / 0,92\\
runs de 5000 & 0,50 / 0,88 & 0,38 / 0,76 & 0,61 / 0,76 & 0,79 / 1,01 & 1,00 / 0,98\\
runs de 1000 & 0,41 / 0,70 & 0,32 / 0,61 & 0,50 / 0,75 & 0,63 / 1,02 & 1,03 / 0,94\\
runs de 63 & 0,34 / 0,50 & 0,26 / 0,42 & 0,41 / 0,56 & 0,50 / 0,65 & 1,00 / 1,01\\
\bottomrule
\end{tabular}
\caption{Tiempo nuevo / tiempo anterior, modo O(1) / modo rápido con un hilo. Menos de 1
es mejora. $n=10^6$ ($250\,000$ \emph{strings}). Las tres últimas filas son tramos
ordenados de claves al azar.}\label{tab:ab}
\end{table}

\noindent Las mejoras grandes están donde antes se perdía: zonas al azar dentro de
datos ordenados (0,45 a 0,81) y runs de claves al azar (0,26 a 0,88, salvo con
comparador en modo rápido, donde no cambia). Las entradas
aleatorias y con repetidos no cambian, salvo los pares (0,76 a 0,81, por la selección
sin salto). Lo que empeoró queda en 10\% o menos, dentro del ruido, con una excepción
de 1,27 en una entrada que tarda menos de un milisegundo. Con \emph{strings} la regla de
ruta y las zonas midieron peor (hasta 1,61), y por eso están desactivadas para tipos
que no se copian byte a byte.

\paragraph{Lo que se probó y no se adoptó.}
\begin{itemize}[nosep]
\item Tamaño del buffer propio (4 a 128\,KB) y de la hoja (32 a 256): las diferencias
  quedaron dentro del ruido. Se mantienen 16\,KB y 128.
\item Desenrollar de a 4 la partición escalar: sin ganancia.
\item Paso de cabeza sin salto dentro de la mezcla con galope: sin ganancia.
\item Selección sin salto también para escalares: 10\% peor en u64, porque ahí gcc ya
  emitía \texttt{cmov}.
\item Mezcla sin salto para runs intercalados que no entran en el buffer: el doble de
  lenta en \emph{organ pipe} y en una sierra regular, donde el predictor acierta.
\end{itemize}

\section{Comparación con otros ordenamientos}

\begin{table}[h]\centering\small
\begin{tabular}{lcccc}
\toprule
 & \multicolumn{3}{c}{SotSort / blitsort} & SotSort / mejor de\\
Entrada & u32 & u64 & clave+índice & Grail, Wiki (32\,B)\\
\midrule
aleatoria & 0,63 & 0,72 & 1,14 & 0,60\\
16 valores & 0,46 & 0,62 & 0,80 & 0,41\\
$\sqrt n$ valores & 0,19 & 0,49 & 0,86 & 0,48\\
Zipf & 0,44 & 0,50 & 0,64 & 0,50\\
ordenada & 1,09 & 1,00 & 2,57 & 0,77\\
invertida & 0,94 & 1,01 & 1,81 & 0,12\\
1\% de ruido & 1,08 & 1,06 & 1,08 & 0,63\\
cola 10\% al azar & 0,67 & 0,75 & 1,50 & 0,82\\
\emph{organ pipe} & 0,58 & 0,70 & 1,13 & 0,69\\
runs de $\sqrt n$ & 0,53 & 0,75 & 1,23 & 1,30\\
runs de 63 & 0,42 & 0,48 & 0,78 & 0,75\\
\bottomrule
\end{tabular}
\caption{Ordenamientos de memoria constante, $n=10^6$: cociente de tiempos, menos de 1
gana SotSort. ``Clave+índice'' es una palabra de 64 bits comparada por su mitad alta, sin
partición vectorial. blitsort no admite registros de 32 bytes.}\label{tab:o1}
\end{table}

\begin{table}[h]\centering\small
\begin{tabular}{lrrrrrr}
\toprule
Entrada (u64) & SotSort O(1) & SotSort rápido & blitsort & fluxsort & glidesort & driftsort\\
\midrule
aleatoria & 23,5 & 21,0 & 32,6 & 27,4 & 30,3 & 26,6\\
Zipf & 13,7 & 13,6 & 27,5 & 16,3 & 28,3 & 13,4\\
ordenada & 0,3 & 0,3 & 0,3 & 0,3 & 0,6 & 0,4\\
invertida & 0,6 & 0,6 & 0,6 & 0,7 & 0,9 & 0,8\\
1\% de ruido & 10,8 & 12,1 & 10,2 & 24,2 & 30,5 & 20,1\\
cola 10\% al azar & 5,6 & 7,6 & 7,4 & 6,9 & 5,1 & 6,6\\
\emph{organ pipe} & 2,8 & 3,6 & 4,0 & 3,2 & 1,9 & 4,9\\
runs de $\sqrt n$ & 25,4 & 22,3 & 33,7 & 26,8 & 12,9 & 44,0\\
runs de 63 & 23,6 & 21,8 & 49,4 & 29,4 & 30,8 & 28,0\\
\bottomrule
\end{tabular}
\caption{Milisegundos, $n=10^6$ enteros de 64 bits, un hilo. ``Rápido'' usa un buffer
completo. Se omiten las filas de 16 y $\sqrt n$ valores, que el modo rápido resuelve con
un conteo.}\label{tab:abs}
\end{table}

\noindent Con enteros de 32 y 64 bits, el modo O(1) queda por debajo de blitsort en 16
de 22 casos (0,19 a 0,75) y a la par en los otros seis (0,94 a 1,09). Con un comparador
que no admite la partición vectorial pierde en 7 de las 11 entradas (1,08 a 2,57; lo
peor, la entrada ordenada) y gana en las otras cuatro (0,64 a 0,86). Contra los ordenamientos con memoria
lineal, el modo rápido gana en aleatoria y pierde con glidesort en \emph{organ pipe},
cola al azar y runs de $\sqrt n$ (1,4 a 1,9 veces en u64). Para $n=10^7$ enteros
aleatorios: 280\,ms el modo rápido, 364 el modo O(1), 362 glidesort, 375 fluxsort, 392
driftsort y 463 blitsort. Memoria adicional con $n=10^7$ (80\,MB de datos): menos de
0,5\,MB en modo O(1) y unos 40\,MB en modo rápido.

\paragraph{Comparaciones y adversarios.} En entrada aleatoria SotSort hace
$1{,}11\lg n!$ comparaciones; todos los medidos están entre 1,07 y 1,29. Contra el
adversario de McIlroy con posiciones congeladas al azar, el máximo fue $1{,}14\,n\lg n$
(0,65 blitsort, 1,96 driftsort, 3,67 ipnsort); ninguno mostró comportamiento
cuadrático. El adversario perezoso del capítulo 1, contra el quicksort solo: con la
línea de la versión 6.1, 20 a 27 comparaciones por elemento para $n$ de 1000 a 8000;
sin ella (\texttt{-DSOT\_V6\_BUG}), 167, 335, 668 y 1335, es decir $n/6$.

\section{Planificación fuera de orden (experimental)}\label{sec:ooo}
Cada rango pendiente del quicksort estable es una tarea, lista apenas termina la
partición que la creó. Los trabajadores toman tareas de una cola acotada; no hay rondas
ni barreras. El buffer compartido era una dependencia falsa: en modo completo la tarea
sobre $[l,r)$ usa $\mathit{aux}[l..r)$. La primera partición se hace en paralelo. La
corrección reutiliza el invariante del quicksort, que habla de un conjunto de
intervalos pendientes (sección 14 de las demostraciones).

\begin{table}[h]\centering\small
\begin{tabular}{lcccccc}
\toprule
 & \multicolumn{2}{c}{u64} & \multicolumn{2}{c}{128 bytes} & \multicolumn{2}{c}{\emph{strings}}\\
Entrada & $10^6$ & $10^7$ & 250 mil & $10^6$ & 250 mil & $10^6$\\
\midrule
aleatoria & 0,85 & 0,83 & 0,72 & 1,04 & 0,59 & 0,78\\
16 valores & 0,67 & 0,77 & 0,69 & 1,09 & 0,76 & 1,10\\
$\sqrt n$ valores & 0,69 & 0,84 & 0,71 & 1,15 & 0,58 & 1,04\\
mitad un solo valor & 0,83 & 0,69 & 0,72 & 0,88 & 0,69 & 0,94\\
mitades desbalanceadas & 0,56 & 0,66 & 0,78 & 0,96 & 0,63 & 0,58\\
1\% de ruido (delegado) & 1,01 & 1,05 & 0,98 & 0,97 & 0,98 & 1,11\\
\emph{organ pipe} (delegado) & 1,17 & 1,28 & 1,00 & 1,06 & 0,92 & 0,94\\
\bottomrule
\end{tabular}
\caption{Tiempo del planificador / tiempo del modo paralelo de v8-r, con 2 hilos. Menos
de 1 gana el planificador. Con 250 mil elementos el modo paralelo no se activa.}\label{tab:ooo}
\end{table}

\noindent Gana donde el modo paralelo no se activa (0,58 a 0,78) y con entradas
desbalanceadas. Con $10^6$ o más queda entre 0,56 y 1,15; con registros de 128 bytes,
entre 0,88 y 1,15. Las entradas con runs se delegan al modo actual (0,92 a 1,28). La
aceleración sobre un hilo fue de 1,27 a 2,43. \textbf{Decisión:} queda como
experimental. El único cambio que los datos justifican es usar tareas cuando
$n<2^{18}$. Solo se midió con 2 núcleos: no hay datos de escalabilidad. La variante con
tareas por bloque y por rotación se descartó sin prototipo, porque su paralelismo
interno es $O(\lg(\lambda/B))$. La idea no es nueva: es quicksort paralelo por tareas
(Cilk, TBB) con un modelo de dependencias como el de StarSs/OmpSs; las citas están
hechas de memoria y hay que verificarlas.

\section{Revisiones r3 a r6: memoria acotada y operaciones contadas}\label{sec:r6}
Esta sección resume lo que cambió entre el 18 y el 19 de noviembre, después de comparar SotSort
con Virtual-Memory Powersort (Moltmann, Nakajima y Wild, ESA 2026) en el arnés de sus autores y
con la misma memoria. Los detalles, demostraciones y pruebas están en
\texttt{docs/IDEA\_particion\_paginada.md}, \texttt{docs/LEMAS\_v8r.md} (observaciones T, P y S)
y el informe técnico, revisión 6.

\paragraph{Cambios.} r3: la densidad de giros se estima con 64 ventanas de 256 pares. r4: mezcla
por páginas cuando los dos runs superan el buffer. r5: runs paginados (esquema de VM-Powersort) y
un presupuesto de memoria que cuenta todo lo que se reserva. r6: (a) la partición estable binaria
cuesta al menos 1{,}5 movimientos por elemento y nivel, más rotaciones que crecen con
$\lg(m/B)$; (b) con buffer acotado y entrada casi aleatoria, un samplesort estable de $k$ vías por
páginas reemplaza la parte alta del quicksort; (c) el quicksort tolera a lo sumo 4 particiones
malas; (d) el grupo de cada elemento se calcula bajando un árbol implícito de separadores sin
saltos.

\paragraph{La pasada por páginas.} Cada grupo tiene una página de salida en el buffer; una página
llena se copia a la primera página del tramo que ya se terminó de leer, y se anota su grupo. Tres
hechos la hacen estable: al volcar la página $f$ ya se leyeron al menos $(f+1)P$ elementos; el
índice de la ranura es el orden de volcado; al final el espacio libre mide exactamente lo que
suman las páginas parciales. Una permutación por ciclos agrupa las páginas y las parciales se
intercalan, de modo que cada grupo queda contiguo. El mecanismo de reutilizar bloques consumidos
viene de Radsort (Clausecker y Schintke, 2026) y el de reordenar bloques de un solo tipo de
Logsort (aphitorite); lo propio es la versión por comparaciones, de arriba hacia abajo, dentro de
un ordenamiento adaptativo.

\paragraph{Resultado.} Con la memoria de VM-Powersort, en las 12 celdas que el arnés puede contar
(\texttt{int} y \texttt{fastarr}, $S$ de 2 a $10^6$), SotSort no hace más comparaciones (salvo
$+0{,}02$ por elemento) ni más movimientos. Con entrada casi aleatoria ($S=2$, $n\approx10^6$,
enteros): 18{,}96 comparaciones y 22{,}10 movimientos por elemento, contra 27{,}76 y 31{,}65.
En tiempo, solo como indicación (dos máquinas virtuales): 0{,}72--0{,}78 de VM-Powersort con
enteros y $S=2$. El adversario de McIlroy, conociendo la semilla, lleva el peor caso medido a unos
34 movimientos por elemento; VM-Powersort, que es un mergesort, no tiene esa debilidad.

\section{Limitaciones}
\begin{itemize}[nosep]
\item Una sola máquina, virtualizada y compartida; sin contadores de hardware; sin ARM.
\item Las constantes de las sondas y de la regla de ruta se ajustaron en esa máquina y
  sobre varias de las mismas entradas con las que después se compara: las filas de datos
  con estructura de las tablas~\ref{tab:o1} y~\ref{tab:abs} son optimistas.
\item Si el comparador lanza una excepción se pueden perder elementos.
\item El samplesort por páginas solo corre con buffer de bloques: en modo $O(1)$ sus tablas no
  entran y sigue el quicksort. Su peor caso está acotado por $O(n\log n)$, pero la constante solo
  se midió contra un adversario.
\item Nada está verificado mecánicamente; la correspondencia entre el pseudocódigo y el
  código se apoya en pruebas diferenciales (19\,506 casos sin fallas, más las variantes
  con AVX2, sin SIMD, clang y los \emph{sanitizers}).
\item Los capítulos 1 y 2 del libro no se pudieron recompilar: se perdieron su archivo
  de estilos y los datos de siete figuras. Se entregan como estaban, con esta adenda al
  final.
\end{itemize}
\end{document}
